\documentclass[
  acmsmall,
  screen,
  natbib=false,
]{acmart}

\usepackage{preamble}

%% BIBLIOGRAPHY

\RequirePackage[
  datamodel=acmdatamodel,
  style=acmauthoryear,
  uniquename=false,
]{biblatex}
\addbibresource{refs.bib}

%% CCS CONCEPTS

\begin{CCSXML}
<ccs2012>
   <concept>
       <concept_id>10003752.10003790.10003795</concept_id>
       <concept_desc>Theory of computation~Constraint and logic programming</concept_desc>
       <concept_significance>500</concept_significance>
       </concept>
   <concept>
       <concept_id>10011007.10011006.10011008.10011009.10011015</concept_id>
       <concept_desc>Software and its engineering~Constraint and logic languages</concept_desc>
       <concept_significance>300</concept_significance>
       </concept>
   <concept>
       <concept_id>10003752.10003753.10010622</concept_id>
       <concept_desc>Theory of computation~Abstract machines</concept_desc>
       <concept_significance>300</concept_significance>
       </concept>
 </ccs2012>
\end{CCSXML}

\ccsdesc[500]{Theory of computation~Constraint and logic programming}
\ccsdesc[300]{Software and its engineering~Constraint and logic languages}
\ccsdesc[300]{Theory of computation~Abstract machines}

\keywords{functional logic programming, narrowing, algebraic effects, implementation}


\begin{document}

\title{\Gweek{}: Implementing Functional Logic Programming from Algebraic Equations}

\author{Samson Main}
\orcid{0009-0002-1580-9032}
\affiliation{%
  \institution{University of Bristol}
  \city{Bristol}
  \country{United Kingdom}
}
\email{ib21231@bristol.ac.uk}

\author{Eddie Jones}
\orcid{0000-0003-1762-5405}
\affiliation{%
  \institution{University of Bristol}
  \city{Bristol}
  \country{United Kingdom}
}
\email{eddie.jones@bristol.ac.uk}

\author{G. A. Kavvos}
\orcid{0000-0001-7953-7975}
\affiliation{%
  \institution{University of Bristol}
  \city{Bristol}
  \country{United Kingdom}
}
\email{alex.kavvos@bristol.ac.uk}

\begin{abstract}

We present \Gweek{}, an implementation of a functional logic programming (FLP)
language whose design is derived directly from the algebraic theory of FLP
effects developed in a companion paper. That paper showed that the
characteristic constructs of FLP---nondeterministic choice, logical variables,
and equational constraints---are algebraic effects, and used a domain-theoretic
semantics based on the lower powerdomain to validate a comprehensive equational
theory and derive an abstract machine. Here we implement a variant of that
machine in Rust. The source language is a small Haskell-like language that is
compiled to a call-by-value core calculus, translated to Call-by-Push-Value
(CBPV), and then executed by a CEK-style abstract machine that explores a
nondeterministic search tree. Because every transition of the machine is
validated by the denotational semantics, the machine is evidently correct, and
the equations of the theory serve directly as the rewrite rules of a peephole
optimiser. We describe a suspension mechanism that defers the evaluation of
\texttt{let} bindings, dramatically improving termination at the cost of some
overhead, and discuss the engineering techniques---arena allocation, persistent
data structures, and four search strategies---that make the implementation
practical.

\end{abstract}

\maketitle

\section{Introduction}
  \label{section:introduction}
  This paper is about implementing \emph{functional logic programming} (FLP)
\parencite{antoy_2010}. FLP extends functional programming with nondeterministic
choice, logical variables, and equational constraints. These features allow a
declarative style of programming in which the programmer describes relationships
between values, and the language searches for solutions. For example, the last
element of a list can be computed by searching for a decomposition:

\begin{haskellcode}
last :: [Nat] -> Nat
last ys = exists xs :: [Nat]. exists x :: Nat. cat xs [x] =:= ys. x.
\end{haskellcode}

\noindent This program introduces logical variables \mintinline{haskell}|xs| and
\mintinline{haskell}|x|, constrains them so that concatenating
\mintinline{haskell}|xs| with the singleton \mintinline{haskell}|[x]| yields
\mintinline{haskell}|ys|, and returns \mintinline{haskell}|x|. The language
discovers the appropriate values automatically through \emph{narrowing}: a
progressive refinement of logical variables driven by the demands of the
computation.

In a companion paper \parencite{jones_2026} we studied the semantics of FLP
through the lens of \emph{algebraic effects} \parencite{plotkin_2003}. The key
observation of that work is that \textbf{the characteristic constructs of FLP
are effects} in the sense of \citeauthor{plotkin_2003}. This led to a core
calculus based on Call-by-Push-Value (CBPV) \parencite{levy_2003}, a
domain-theoretic denotational semantics based on the lower powerdomain, and a
comprehensive equational theory that validates a range of program
transformations. Crucially, the denotational semantics gives rise---via a
\emph{representation theorem}---to an abstract machine whose transitions are
justified by the semantics. The equational theory, in turn, provides a catalogue
of valid program equations.

The present paper asks: \emph{can this theory serve as a practical guide for
implementation?} We answer in the affirmative by presenting
\Gweek{},\!\footnote{%
  Available at \url{https://github.com/xxx/gweek}. The name is a portmanteau of
  the Welsh \emph{gweek} (`search') and the characteristic exclamation of the
  language user upon discovering a particularly elegant solution.}
a small but complete FLP language whose implementation is derived directly from
the theory. \Gweek{} implements a variant of the abstract machine of
\parencite{jones_2026}, and because every transition of the machine is validated
by the denotational semantics, the machine is \emph{evidently correct}: its
correctness is not established by ad hoc reasoning, but follows directly from
the semantic model.

Three aspects of this connection are particularly illuminating.

First, the equational theory of \parencite{jones_2026} serves directly as the
rewrite rules of a \emph{peephole optimiser}
(\cref{section:optimisation}). Each rewrite rule applied by the optimiser is a
validated equation of the theory; its soundness follows from the full abstraction
theorem of the denotational semantics.

Second, the CBPV distinction between values and computations leads naturally to a
\emph{suspension mechanism} for \texttt{let} bindings
(\cref{section:suspensions}). This mechanism defers the evaluation of bound
computations, dramatically improving termination by allowing constraints to prune
the search tree before the full nondeterministic tree is materialised. The
equational theory illuminates precisely when and why this matters.

Third, the machine's treatment of logical variables and narrowing
(\cref{section:machine}) is a direct realisation of the abstract machine of
\parencite{jones_2026}. A CBPV term gives rise to a nondeterministic search
tree; the implementation explores this tree using one of four strategies. The
narrowing strategy---which refines a logical variable only when a case expression
demands its value---corresponds to \emph{weak head narrowing}, and inherits the
completeness guarantee of \parencite{jones_2026}.

\paragraph{Contributions}

The contributions of this paper are:
\begin{itemize}
  \item A description of \Gweek{}, a typed higher-order FLP language. We present
    its surface syntax, a CBV core calculus, and its translation to CBPV
    (\cref{section:source,section:ir}).
  \item A presentation of the equational theory of CBPV with FLP effects, its
    restriction to CBV equations, and an example proof showing how CBPV equations
    validate CBV transformations (\cref{section:ir}).
  \item A CEK-style abstract machine for CBPV that explores nondeterministic
    search trees, with four search strategies (\cref{section:machine}).
  \item A description of the Rust implementation, including arena allocation,
    persistent data structures, and the compilation pipeline
    (\cref{section:implementation}).
  \item A peephole optimiser whose rewrite rules are drawn from the equational
    theory (\cref{section:optimisation}).
  \item A suspension mechanism that improves termination by deferring
    evaluation, with an analysis of its interaction with nondeterminism
    (\cref{section:suspensions}).
\end{itemize}


\section{The Source Language}
  \label{section:source}
  We begin by describing the \Gweek{} language through examples, and then distill
its essential features into a core calculus.

\subsection{The Language by Example}

A \Gweek{} program consists of a sequence of typed function definitions followed
by a main expression. The syntax is designed to resemble a small subset of
Haskell. Consider the following program, which concatenates two lists:
%
\begin{haskellcode}
cat :: [Nat] -> [Nat] -> [Nat]
cat xs ys = case xs of
      [] -> ys
    | (x:xs) -> x : (cat xs ys).
\end{haskellcode}
%
\noindent This is a standard recursive definition: the type signature declares
that \mintinline{haskell}|cat| takes two lists of natural numbers and returns a
list of natural numbers; the definition proceeds by pattern matching on the first
argument.

What distinguishes \Gweek{} from a purely functional language are three
additional primitives. First, \emph{nondeterministic choice}
(\mintinline{haskell}|<>|) allows a function to return multiple results. The
function \mintinline{haskell}|insert| below nondeterministically inserts an
element at every position in a list:
%
\begin{haskellcode}
insert :: Nat -> [Nat] -> [Nat]
insert x xs = case xs of
      [] -> [x]
    | (z:zs) -> ((x : z : zs) <> (z : insert x zs)).
\end{haskellcode}
%
\noindent The expression \mintinline{haskell}|(x : z : zs) <> (z : insert x zs)|
evaluates to \emph{both} alternatives: one places \mintinline{haskell}|x| at
the front, the other recurses deeper into the list. The nullary version,
\mintinline{haskell}|fail|, produces no results at all.

Second, \emph{logical variables} (\mintinline{haskell}|exists|) introduce fresh
unknowns that are progressively refined by the machine:
%
\begin{haskellcode}
last :: [Nat] -> Nat
last xs = exists ys :: [Nat]. exists y :: Nat.
    cat ys [y] =:= xs. y.
\end{haskellcode}
%
\noindent Here \mintinline{haskell}|ys| and \mintinline{haskell}|y| are
logical variables. The third primitive, the \emph{equational constraint}
(\mintinline{haskell}|=:=|), attempts to unify two values. If unification
succeeds, execution continues under the resulting substitution; if it fails, the
current branch is discarded. The program thus computes the last element of
\mintinline{haskell}|xs| by searching for a decomposition into an initial
segment \mintinline{haskell}|ys| and a final element \mintinline{haskell}|y|.

The interaction between logical variables and pattern matching gives rise to
\emph{narrowing}. When the machine encounters a case expression whose scrutinee
is a logical variable, it nondeterministically splits into branches, one for each
possible head constructor of the scrutinee's type. In the example above, calling
\mintinline{haskell}|cat ys [y]| pattern-matches on \mintinline{haskell}|ys|,
which triggers narrowing: the machine splits into a branch where
$\mathit{ys} = \texttt{[]}$ and a branch where $\mathit{ys} = h : t$ for fresh
variables $h$ and $t$, and continues recursively.

\subsection{The CBV Core Calculus}

\begin{figure*}
  \begin{small}
\begin{minipage}{0.28\textwidth}
  {\color{CBVGreen}\fbox{\color{black}\text{types}}}
  \[
    \begin{array}{rcl}
           A     & \Coloneqq\ & \UnitT \mid \NatT \\
                 &            & A_1 + A_2 \\
                 &            & A_1 \times A_2 \\
                 &            & \ListT{A} \\
                 &            & A_1 \to A_2
    \end{array}
  \]
  {\color{CBVGreen}\fbox{\color{black}\text{first-order types}}}
  \[
    \begin{array}{rcl}
      \sigma    & \Coloneqq\ & \UnitT \mid \NatT \\
                &            & \sigma_1 + \sigma_2 \\
                &            & \sigma_1 \times \sigma_2 \\
                &            & \ListT{\sigma}
    \end{array}
  \]
\end{minipage}
\begin{minipage}{0.70\textwidth}
  {\color{CBVGreen}\fbox{\color{black}\text{expressions}}} \; $\CBV{e}$
  \[
    \begin{array}{rcll}
      \CBV{e}   & \Coloneqq\ & \CBV{x}                                          & \text{variable} \\
                &            & \CBV{\UnitV} \mid \CBV{\Zero} \mid \CBV{\Succ{e}} & \text{unit, zero, successor} \\
                &            & \CBV{\VTuple{e_1}{e_2}} \mid \CBV{\In[i]{e}}      & \text{pair, injection} \\
                &            & \CBV{\Nil} \mid \CBV{\Cons{e_1}{e_2}}             & \text{nil, cons} \\
                &            & \CBV{\CBVLam{x}[A]{e}}                            & \text{abstraction} \\
                &            & \CBV{\CBVApp{e_1}{e_2}}                           & \text{application} \\
                &            & \CBV{\CBVLet{x}{e_1}{e_2}}                        & \text{binding} \\
                &            & \CBV{\CBVIfz{e}{e_0}{n}{e_1}}                     & \text{nat case} \\
                &            & \CBV{\CBVMatch{e}{e_0}{x}{xs}{e_1}}              & \text{list case} \\
                &            & \CBV{\CBVSplit{e}{x_1}{x_2}{e'}}                 & \text{pair elim.} \\
                &            & \CBV{\CBVCase{e}{x}{e_1}{y}{e_2}}               & \text{sum case} \\
                &            & \CBV{\CBVFix{f}[A]{e}}                            & \text{recursion} \\[0.3ex]
                &            & \CBV{\Fail} \mid \CBV{\Choice{}{e_1}{e_2}}        & \text{failure, choice} \\
                &            & \CBV{\Exists{x}[\sigma]{e}}                       & \text{fresh variable} \\
                &            & \CBV{\CBVEquate{e_1}{e_2}{e}}                     & \text{unification}
    \end{array}
  \]
\end{minipage}
\end{small}
 \\[1ex]
  \begin{footnotesize}
\begin{mathpar}
  \inferH{Var}{
    \strut
  }{
    \IsCBV[\Gamma, x : A, \Gamma']{x}{A}
  }
  \and
  \inferH{Unit}{
  }{
    \IsCBV{\UnitV}{\UnitT}
  }
  \and
  \inferH{Zero}{
    \strut
  }{
    \IsCBV{\Zero}{\NatT}
  }
  \and
  \inferH{Succ}{
    \IsCBV{E}{\NatT}
  }{
    \IsCBV{\Succ{E}}{\NatT}
  }
  \and
  \inferH{Pair}{
    \IsCBV{E_1}{A_1} \\
    \IsCBV{E_2}{A_2}
  }{
    \IsCBV{\VTuple{E_1}{E_2}}{A_1 \times A_2}
  }
  \and
  \inferH{In}{
    \IsCBV{E}{A_i}
  }{
    \IsCBV{\In[i]{E}}{A_1 + A_2}
  }
  \and
  \inferH{Nil}{
    \strut
  }{
    \IsCBV{\Nil}{\ListT{A}}
  }
  \and
  \inferH{Cons}{
    \IsCBV{E_1}{A} \\
    \IsCBV{E_2}{\ListT{A}}
  }{
    \IsCBV{\Cons{E_1}{E_2}}{\ListT{A}}
  }
  \and
  \inferH{Lam}{
    \IsCBV[\Gamma, \DeclVar{x}{A_1}]{E}{A_2}
  }{
    \IsCBV{\CBVLam{x}[A_1]{E}}{A_1 \to A_2}
  }
  \and
  \inferH{App}{
    \IsCBV{E_1}{A_1 \to A_2} \\
    \IsCBV{E_2}{A_1}
  }{
    \IsCBV{\CBVApp{E_1}{E_2}}{A_2}
  }
  \and
  \inferH{Let}{
    \IsCBV{E_1}{A_1} \\
    \IsCBV[\Gamma, \DeclVar{x}{A_1}]{E_2}{A_2}
  }{
    \IsCBV{\CBVLet{x}{E_1}{E_2}}{A_2}
  }
  \and
  \inferH{Rec}{
    \IsCBV[\Gamma, f : A]{E}{A}
  }{
    \IsCBV{\CBVFix{f}[A]{E}}{A}
  }
  \and
  \inferH{Ifz}{
    \IsCBV{E}{\NatT} \\
    \IsCBV{E_0}{A} \\
    \IsCBV[\Gamma, n : \NatT]{E_1}{A}
  }{
    \IsCBV{\CBVIfz{E}{E_0}{n}{E_1}}{A}
  }
  \and
  \inferH{Match}{
    \IsCBV{E}{\ListT{A'}} \\
    \IsCBV{E_0}{A} \\
    \IsCBV[\Gamma, x : A', xs : \ListT{A'}]{E_1}{A}
  }{
    \IsCBV{\CBVMatch{E}{E_0}{x}{xs}{E_1}}{A}
  }
  \and
  \inferH{Split}{
    \IsCBV{E}{A_1 \times A_2} \\
    \IsCBV[\Gamma, x_1 : A_1, x_2 : A_2]{E'}{A}
  }{
    \IsCBV{\CBVSplit{E}{x_1}{x_2}{E'}}{A}
  }
  \and
  \inferH{Case}{
    \IsCBV{E}{A_1 + A_2} \\
    \IsCBV[\Gamma, x : A_1]{E_1}{A} \\
    \IsCBV[\Gamma, y : A_2]{E_2}{A}
  }{
    \IsCBV{\CBVCase{E}{x}{E_1}{y}{E_2}}{A}
  }
  \and
  \inferH{Fail}{
    \strut
  }{
    \IsCBV{\Fail}{A}
  }
  \and
  \inferH{Choice}{
    \IsCBV{E_1}{A} \\
    \IsCBV{E_2}{A}
  }{
    \IsCBV{\Choice{}{E_1}{E_2}}{A}
  }
  \and
  \inferH{Exists}{
    \IsCBV[\Gamma, x : \sigma]{E}{A}
  }{
    \IsCBV{\Exists{x}[\sigma]{E}}{A}
  }
  \and
  \inferH{Equate}{
    \IsCBV{E_1}{\sigma} \\
    \IsCBV{E_2}{\sigma} \\
    \IsCBV{E}{A}
  }{
    \IsCBV{\CBVEquate{E_1}{E_2}{E}}{A}
  }
\end{mathpar}
\end{footnotesize}

  \caption{The CBV Core Calculus: Syntax and Type System
    {\color{CBVGreen}(green)}. The last four rules (\ruleref{Fail},
    \ruleref{Choice}, \ruleref{Exists}, \ruleref{Equate}) are the FLP effects.}
  \label{figure:cbv}
\end{figure*}

We now distill the source language into a core calculus (\cref{figure:cbv}). The
calculus is a standard call-by-value $\lambda$-calculus extended with algebraic
data types, recursion, and the three FLP effects. There is a single syntactic
category of expressions ($\CBV{e}$) and a single typing judgment $\IsCBV{e}{A}$.

\paragraph{Types}
The type system includes natural numbers ($\NatT$), lists ($\ListT{A}$),
products ($A_1 \times A_2$), sums ($A_1 + A_2$), and function types ($A_1 \to
A_2$). The \emph{first-order types} $\sigma$ are the subset without function
types; these are the types over which logical variables may range. The
restriction to first-order types is both theoretically
motivated---first-order types admit the simple domain-theoretic semantics of
\parencite{jones_2026}---and practically necessary, as equational constraints on
function types would require higher-order unification, which is undecidable
\parencite{huet_1973}. The surface language also provides booleans and
\mintinline{haskell}|if|/\mintinline{haskell}|then|/\mintinline{haskell}|else|,
which are desugared into sums and case analysis.

\paragraph{Expressions}
Expressions include data constructors (\textsf{zero}, \textsf{succ},
\textsf{nil}, \textsf{cons}, pairs, injections), $\lambda$-abstractions,
application, let-bindings, pattern matching (decomposed into eliminators for each
type), and recursion ($\textsf{rec}$). The FLP primitives---failure
($\textsf{fail}$), nondeterministic choice ($\talloblong$), logical variables
($\exists x : \sigma.\, e$), and equational constraints ($e_1
\mathbin{\text{\textdblhyphen}} e_2;\, e$)---complete the calculus.

\paragraph{Recursion and function definitions}
Each top-level function definition is compiled to a recursion
$\CBV{\CBVFix{f}[A]{e}}$, which binds $f$ to itself in the body $\CBV{e}$.
Function definitions may be mutually recursive; the compiler topologically sorts
them and compiles each group into nested fixed points.


\section{Intermediate Representation}
  \label{section:ir}
  The CBV core calculus is translated into Call-by-Push-Value (CBPV)
\parencite{levy_2003} before execution. This intermediate representation is not
merely a compilation target: it is the calculus in which the equational theory of
\parencite{jones_2026} is formulated, and whose equations justify both the
transitions of the abstract machine and the rewrite rules of the peephole
optimiser. In this section we describe the translation, present the equational
theory, and show how it validates program transformations.

\subsection{Translation to CBPV}
\label{subsection:translation}

\begin{figure*}
  \begin{footnotesize}
\begin{minipage}{0.28\textwidth}
{\color{CBVGreen}\fbox{\color{black}\text{type translation}}}
\; $\Transl{A}$
\begin{align*}
    \Transl{\UnitT} &= \UnitT \\
    \Transl{\NatT} &= \NatT \\
    \Transl{A_1 + A_2} &= \Transl{A_1} + \Transl{A_2} \\
    \Transl{A_1 \times A_2} &= \Transl{A_1} \times \Transl{A_2} \\
    \Transl{\ListT{A}} &= \ListT{\Transl{A}} \\
    \Transl{A_1 \to A_2} &= \ThunkT*{\Transl{A_1} \to \ReturnerT*{\Transl{A_2}}}
\end{align*}
\end{minipage}
\begin{minipage}{0.70\textwidth}
{\color{CBVGreen}\fbox{\color{black}\text{term translation}}}
\; $\Transl{\CBV{e}} = \CTerm{M}$
\begin{align*}
    \Transl{\CBV{x}} &= \CTerm{\Return{x}}
    \\
    \Transl{\CBV{\UnitV}} &= \CTerm{\Return{\UnitV}}
    \qquad
    \Transl{\CBV{\Zero}} = \CTerm{\Return{\Zero}}
    \\
    \Transl{\CBV{\Succ{e}}} &= \CTerm{\CTo{\Transl*{\CBV{e}}}{v}{\Return{\Succ{v}}}}
    \\
    \Transl{\CBV{\VTuple{e_1}{e_2}}} &= \CTerm{\CTo{\Transl*{\CBV{e_1}}}{v_1}{\CTo*{\Transl*{\CBV{e_2}}}{v_2}{\Return{\VTuple{v_1}{v_2}}}}}
    \\
    \Transl{\CBV{\Nil}} &= \CTerm{\Return{\Nil}}
    \\
    \Transl{\CBV{\Cons{e_1}{e_2}}} &= \CTerm{\CTo{\Transl*{\CBV{e_1}}}{v}{\CTo*{\Transl*{\CBV{e_2}}}{w}{\Return{\Cons{v}{w}}}}}
    \\
    \Transl{\CBV{\CBVLam{x}{e}}} &= \CTerm{\Return*{\Thunk{\Lam{x}{\Transl{\CBV{e}}}}}}
    \\
    \Transl{\CBV{\CBVApp{e_1}{e_2}}} &= \CTerm{\CTo{\Transl*{\CBV{e_1}}}{f}{\CTo*{\Transl*{\CBV{e_2}}}{x}{\CApp{\Force{f}}{x}}}}
    \\
    \Transl{\CBV{\CBVLet{x}{e_1}{e_2}}} &= \CTerm{\CTo{\Transl*{\CBV{e_1}}}{x}{\Transl*{\CBV{e_2}}}}
    \\
    \Transl{\CBV{\CBVIfz{e}{\ldots}{\ldots}{\ldots}}} &= \CTerm{\CTo{\Transl*{\CBV{e}}}{v}{\CIfz{v}{\Transl*{\CBV{e_0}}}{n}{\Transl*{\CBV{e_1}}}}}
    \\
    \Transl{\CBV{\Fail}} &= \CTerm{\Fail}
    \qquad
    \Transl{\CBV{\Choice{}{e_1}{e_2}}} = \CTerm{\Choice{}{\Transl*{\CBV{e_1}}}{\Transl*{\CBV{e_2}}}}
    \\
    \Transl{\CBV{\Exists{x}[\sigma]{e}}} &= \CTerm{\Exists{x}[\sigma]{\Transl*{\CBV{e}}}}
    \\
    \Transl{\CBV{\CBVEquate{e_1}{e_2}{e}}} &= \CTerm{\CTo{\Transl*{\CBV{e_1}}}{v}{\CTo*{\Transl*{\CBV{e_2}}}{w}{\Equate{v}{w}{\Transl*{\CBV{e}}}}}}
\end{align*}
\end{minipage}
\end{footnotesize}

  \caption{Translation from {\color{CBVGreen}CBV} to
    {\color{ACMMatch1}CBPV} {\color{ACMMatch2}terms}. Each CBV expression
    $\CBV{e}$ of type $A$ translates to a CBPV computation $\CTerm{M}$ of
    type $\ReturnerT*{\Transl{A}}$. Constructor arguments are sequenced via
    $\textsf{to}$, making the evaluation order explicit.}
  \label{figure:translation}
\end{figure*}

The translation (\cref{figure:translation}) follows the standard embedding of
call-by-value into CBPV \parencite[\S 4]{levy_2003}, extended with the FLP
effects. Each CBV expression $\CBV{e}$ of type $A$ becomes a CBPV computation of
type $\ReturnerT*{\Transl{A}}$.

The key type translation is for functions: a CBV type $A_1 \to A_2$ becomes the
CBPV value type $\ThunkT*{\Transl{A_1} \to \ReturnerT*{\Transl{A_2}}}$. The
outer $\ThunkT{(-)}$ makes the function a storable value; the inner
$\ReturnerT{(-)}$ marks the result as an effectful computation. All other types
translate homomorphically.

The term translation makes evaluation order explicit. A variable $\CBV{x}$
becomes $\CTerm{\Return{x}}$: a trivially completed computation. A constructor
with subexpressions, such as $\CBV{\Succ{e}}$, must first evaluate its argument:
$\CTerm{\CTo{\Transl*{\CBV{e}}}{v}{\Return{\Succ{v}}}}$. Application
$\CBV{\CBVApp{e_1}{e_2}}$ evaluates both subexpressions, forces the resulting
thunk, and applies it. A $\lambda$-abstraction $\CBV{\CBVLam{x}{e}}$ becomes
$\CTerm{\Return*{\Thunk{\Lam{x}{\Transl{\CBV{e}}}}}}$. The FLP effects
translate directly: $\CBV{\Fail}$ becomes $\CTerm{\Fail}$,
$\CBV{\Choice{}{e_1}{e_2}}$ becomes $\CTerm{\Choice{}{M_1}{M_2}}$, and so on.

\paragraph{Why CBPV?}
The value/computation distinction is essential for three reasons. First, it makes
the sequencing of effects explicit: every point at which an effect may occur is
visible in the CBPV term, which is critical for reasoning about nondeterminism.
Second, CBPV provides the framework in which the equational theory of
\parencite{jones_2026} is formulated; without it, the equations that justify our
optimiser would not be available. Third, the CBPV structure maps directly onto
the abstract machine: values are data, computations are processes, and the
machine's stack tracks pending continuations.

\subsection{The Equational Theory}
\label{subsection:equations}

\begin{figure*}
  \begin{footnotesize}
\fbox{\textbf{$\beta$-laws}} \\[0.5ex]
\begin{minipage}{0.48\textwidth}
\begin{align*}
  \CTerm{\CTo*{\Return{V}}{x}{M}}
    &= \CTerm{\CSubst{M}{V}{x}} \\
  \CTerm{\CApp*{\Lam{x}{M}}{V}}
    &= \CTerm{\CSubst{M}{V}{x}} \\
  \CTerm{\Force*{\Thunk{M}}}
    &= \CTerm{M}
\end{align*}
\end{minipage}
\begin{minipage}{0.48\textwidth}
\begin{align*}
  \CTerm{\CIfz{\Zero}{M}{n}{N}}
    &= \CTerm{M} \\
  \CTerm{\CIfz{\Succ{V}}{M}{n}{N}}
    &= \CTerm{\CSubst{N}{V}{n}} \\
  \CTerm{\CMatch{\Cons{V}{W}}{M}{x}{xs}{N}}
    &= \CTerm{\SubstTwo{N}{V}{W}{x}{xs}}
\end{align*}
\end{minipage}

\bigskip

\fbox{\textbf{$\eta$-laws}} \\[0.5ex]
\begin{minipage}{0.48\textwidth}
\begin{align*}
  \CTerm{M}
    &= \CTerm{\CTo{M}{x}{\Return{x}}} \\
  \CTerm{M}
    &= \CTerm{\Lam{x}{\CApp{M}{x}}}
\end{align*}
\end{minipage}
\begin{minipage}{0.48\textwidth}
\begin{align*}
  \VTerm{V}
    &= \VTerm{\Thunk*{\Force{V}}} \\
  \CTerm{\CSubst{M}{V}{n}}
    &= \CTerm{\CIfz{V}{\CSubst{M}{\Zero}{n}}{m}{\CSubst{M}{\Succ{m}}{n}}}
\end{align*}
\end{minipage}

\bigskip

\fbox{\textbf{Sequencing laws}} \\[0.5ex]
\begin{align*}
  \CTerm{\CTo*{\CTo{M}{x}{N}}{y}{P}}
    &= \CTerm{\CTo{M}{x}{(\CTo{N}{y}{P})}}
  & \CTerm{\CApp*{\CTo{M}{x}{N}}{V}}
    &= \CTerm{\CTo{M}{x}{\CApp{N}{V}}}
\end{align*}

\bigskip

\fbox{\textbf{Effect equations}} \\[0.5ex]
\begin{minipage}{0.48\textwidth}
\begin{align*}
  \CTerm{\Choice{\Fail}{M}} &= \CTerm{M} \\
  \CTerm{\Choice{M}{\Fail}} &= \CTerm{M} \\
  \CTerm{\Choice{M}{N}} &= \CTerm{\Choice{N}{M}} \\
  \CTerm{\Choice{M}{M}} &= \CTerm{M}
\end{align*}
\end{minipage}
\begin{minipage}{0.48\textwidth}
\begin{align*}
  \CTerm{\Equate{V}{V}{M}} &= \CTerm{M} \\
  \CTerm{\CTo{M}{x}{\Fail}} &= \CTerm{\Fail} \\
  \CTerm{\Exists{x}[\sigma]{\Fail}} &= \CTerm{\Fail} \\
  \CTerm{\Choice{M_1}{\DelimPrn{\Choice{M_2}{M_3}}}}
    &= \CTerm{\Choice{\DelimPrn{\Choice{M_1}{M_2}}}{M_3}}
\end{align*}
\end{minipage}

\bigskip

\fbox{\textbf{Pull-through laws}} \\[0.5ex]
\begin{align*}
  \CTerm{\CTo*{\Choice{M_1}{M_2}}{x}{N}}
    &= \CTerm{\Choice{(\CTo{M_1}{x}{N})}{(\CTo{M_2}{x}{N})}}
  \\
  \CTerm{\CTo*{\Exists{z}[\sigma]{M}}{x}{N}}
    &= \CTerm{\Exists*{z}[\sigma]{\CTo{M}{x}{N}}}
  \\
  \CTerm{\CTo*{\Equate{V}{W}{M}}{x}{N}}
    &= \CTerm{\Equate{V}{W}{\CTo{M}{x}{N}}}
\end{align*}

\bigskip

\fbox{\textbf{$\lambda$-laws}} \\[0.5ex]
\begin{align*}
  \CTerm{\Lam{x}{\Fail}}
    &= \CTerm{\Fail}
  &
  \CTerm{\Lam*{x}{\Choice{M}{N}}}
    &= \CTerm{\Choice{\Lam{x}{M}}{\Lam{x}{N}}}
  \\
  \CTerm{\Lam*{x}{\Exists{z}[\sigma]{M}}}
    &= \CTerm{\Exists*{z}[\sigma]{\Lam{x}{M}}}
  &
  \CTerm{\Lam*{x}{\Equate{V}{W}{M}}}
    &= \CTerm{\Equate{V}{W}{\Lam{x}{M}}}
\end{align*}

\fbox{\textbf{Parameter equations} (example: $\NatT$)} \\[0.5ex]
\begin{align*}
  \CTerm{\Exists{n}[\NatT]{M}}
    &= \CTerm{\Choice{\CSubst{M}{\Zero}{n}}{\Exists{m}[\NatT]{\CSubst{M}{\Succ{m}}{n}}}}
\end{align*}

\fbox{\textbf{Narrowing equations} (examples)} \\[0.5ex]
\begin{align*}
  \CTerm{\Equate{\Succ{V}}{\Succ{W}}{M}}
    &= \CTerm{\Equate{V}{W}{M}}
  &
  \CTerm{\Equate{\Succ{V}}{\Zero}{M}}
    &= \CTerm{\Fail}
\end{align*}
\end{footnotesize}

  \caption{Selected equations of the CBPV equational theory, validated by the
    denotational semantics of \parencite{jones_2026}. All equations are between
    well-typed CBPV terms.}
  \label{figure:equations}
\end{figure*}

The denotational semantics of \parencite{jones_2026} validates a rich equational
theory for CBPV with FLP effects. These equations serve two purposes: they
justify the transitions of the abstract machine (\cref{section:machine}), and
they provide the rewrite rules for the peephole optimiser
(\cref{section:optimisation}). We collect the most important equations in
\cref{figure:equations} and discuss them briefly.

\paragraph{Structural equations}
The $\beta$-laws express the computational content of each eliminator: case
analysis on a known constructor reduces, application of a $\lambda$-abstraction
substitutes, forcing a thunk unwraps it, and the bind-return law
$\CTerm{\CTo{\Return{V}}{x}{M}} = \CTerm{\CSubst{M}{V}{x}}$ expresses that
returning a value and immediately binding it is the same as substitution. The
$\eta$-laws express the universality of each type's eliminator. The sequencing
laws express associativity of bind and its interaction with application.

\paragraph{Effect equations}
The algebraic structure of the FLP effects gives rise to a family of equations.
Failure is the identity for choice: $\CTerm{\Choice{\Fail}{M}} = \CTerm{M}$.
Choice is associative, commutative, and idempotent:
$\CTerm{\Choice{M}{M}} = \CTerm{M}$. Existential quantifiers commute with each
other and distribute over choice. Equational constraints are reflexive
($\CTerm{\Equate{V}{V}{M}} = \CTerm{M}$), commute with each other, and
distribute over choice. The \emph{dead-end} law states that
$\CTerm{\CTo{M}{x}{\Fail}} = \CTerm{\Fail}$: if the continuation always fails,
the preceding computation is irrelevant.

\paragraph{Pull-through laws}
The \emph{pull-through} laws describe how effects interact with sequencing. If
the bound computation branches, the continuation runs in each branch:
$\CTerm{\CTo*{\Choice{M_1}{M_2}}{x}{N}} =
\CTerm{\Choice{(\CTo{M_1}{x}{N})}{(\CTo{M_2}{x}{N})}}$. Existential
quantifiers and equational constraints can be pulled through a bind:
$\CTerm{\CTo*{\Exists{z}[\sigma]{M}}{x}{N}} =
\CTerm{\Exists*{z}[\sigma]{\CTo{M}{x}{N}}}$. These laws are used extensively by
the optimiser to simplify machine terms.

\paragraph{Parameter and narrowing equations}
The \emph{parameter equations} characterise the behaviour of $\exists$ on each
first-order type. For example, $\CTerm{\Exists{n}[\NatT]{M}} =
\CTerm{\Choice{\CSubst{M}{\Zero}{n}}{\Exists{m}[\NatT]{\CSubst{M}{\Succ{m}}{n}}}}$:
introducing a logical variable of type $\NatT$ is the same as choosing between
zero and the successor of a fresh variable. The \emph{narrowing equations} govern
unification: matching constructors are peeled off
($\CTerm{\Equate{\Succ{V}}{\Succ{W}}{M}} = \CTerm{\Equate{V}{W}{M}}$) and
mismatched constructors cause failure
($\CTerm{\Equate{\Succ{V}}{\Zero}{M}} = \CTerm{\Fail}$).

\subsection{CBV Equations}
\label{subsection:cbv-equations}

\begin{figure}
\begin{small}
\begin{align}
  \CBV{\CBVLet{x}{(\CBVLet{y}{e_1}{e_2})}{e_3}}
    &= \CBV{\CBVLet{y}{e_1}{(\CBVLet{x}{e_2}{e_3})}}
  \label{eq:cbv-assoc} \tag{Assoc} \\
  \CBV{\CBVLet{x}{e}{x}}
    &= \CBV{e}
  \label{eq:cbv-eta} \tag{$\eta$-Let} \\
  \CBV{\CBVLet{x}{\Fail}{e}}
    &= \CBV{\Fail}
  \label{eq:cbv-fail} \tag{Dead-End} \\
  \CBV{\CBVLet{x}{(\Choice{}{e_1}{e_2})}{e}}
    &= \CBV{\Choice{}{(\CBVLet{x}{e_1}{e})}{(\CBVLet{x}{e_2}{e})}}
  \label{eq:cbv-choice} \tag{$\talloblong$-Let} \\
  \CBV{\Choice{}{e}{\Fail}}
    &= \CBV{e}
  \label{eq:cbv-fail-id} \tag{$\talloblong$-Id} \\
  \CBV{\Choice{}{e}{e}}
    &= \CBV{e}
  \label{eq:cbv-idem} \tag{$\talloblong$-Idem}
\end{align}
\end{small}
\caption{Selected CBV equations, valid by translation to CBPV.}
\label{figure:cbv-equations}
\end{figure}

The CBPV equations induce equations on the CBV calculus. If $\Transl{\CBV{e_1}}$
and $\Transl{\CBV{e_2}}$ are equal as CBPV terms (using the equations of
\cref{figure:equations}), then $\CBV{e_1}$ and $\CBV{e_2}$ are
observationally equivalent. \Cref{figure:cbv-equations} lists several such
equations. These are the equations that a programmer can rely on when reasoning
about \Gweek{} programs.

\paragraph{Example proof}
We demonstrate the method by proving \eqref{eq:cbv-assoc}: the associativity of
let-bindings. Translating both sides to CBPV, the left-hand side becomes
$\CTerm{\CTo{(\CTo{\Transl*{\CBV{e_1}}}{y}{\Transl*{\CBV{e_2}}})}{x}{\Transl*{\CBV{e_3}}}}$
and the right-hand side becomes
$\CTerm{\CTo{\Transl*{\CBV{e_1}}}{y}{(\CTo{\Transl*{\CBV{e_2}}}{x}{\Transl*{\CBV{e_3}}})}}$.
These are equal by the CBPV sequencing law
$\CTerm{\CTo*{\CTo{M}{x}{N}}{y}{P}} =
\CTerm{\CTo{M}{x}{(\CTo{N}{y}{P})}}$ from \cref{figure:equations}.

\paragraph{A plausible but invalid equation}
Not every equation that appears plausible is valid. Consider:
%
\[
  \CBV{\CBVLet{x}{e_1}{\Choice{}{e_2}{e_3}}}
  \stackrel{?}{=}
  \CBV{\Choice{}{(\CBVLet{x}{e_1}{e_2})}{(\CBVLet{x}{e_1}{e_3})}}
\]
%
This `right-distribution' of let over choice would state that binding over a
choice is the same as choosing between bindings. In a deterministic language this
would be harmless, but in FLP it is \emph{not} valid: if $\CBV{e_1}$ is
nondeterministic (e.g.\ $\CBV{e_1} = \CBV{\Choice{}{\Zero}{\Succ{\Zero}}}$),
the left-hand side evaluates $\CBV{e_1}$ once and branches on the continuation,
while the right-hand side evaluates $\CBV{e_1}$ twice, potentially doubling the
search tree. The CBPV translation makes this clear: the left-hand side is
$\CTerm{\CTo{M}{x}{\Choice{N_1}{N_2}}}$, which has a single copy of $\CTerm{M}$
above the choice, while the right-hand side is
$\CTerm{\Choice{(\CTo{M}{x}{N_1})}{(\CTo{M}{x}{N_2})}}$, which duplicates
$\CTerm{M}$. These are not equal: the commutativity law (which allows reordering
of independent binds) does not apply because $\CBV{e_2}$ and $\CBV{e_3}$ may
depend on $x$. The idempotence of choice
($\CTerm{\Choice{M}{M}} = \CTerm{M}$) does not help either, because the
duplicated copies may produce different \emph{numbers} of results, not just
different results.


\section{The Abstract Machine}
  \label{section:machine}
  \begin{figure*}
  \begin{minipage}{0.49\textwidth}
  \[
  \begin{array}{lrcl}
    \text{stacks}
        & \Stk{K} & \Coloneqq\  & \Stk{\StkEmpty}         \\
        &         &             & \Stk{\StkKont{x}{N}{K}} \\
        &         &             & \Stk{\StkVal{V}{K}}     \\[0.8ex]
    \text{logic ctx.}
        & \Gamma  & \Coloneqq\  & \\
        &         &             & \Gamma, x : \sigma      \\[0.8ex]
    \text{susp.\ env.}
        & \Susp{S} & \Coloneqq & \SuspEmpty{} \\
        &          &           & \AddSusp{S}{x}{M}
  \end{array}
  \]
\end{minipage}
\begin{minipage}{0.49\textwidth}
  \[
  \begin{array}{lrcl}
    \text{basic machine}
        & \Machine{L} & \Coloneqq & \SLMachine{M}{K} \\[0.8ex]
    \text{configurations}
        & \Conf{C} & \Coloneqq & \FailC{}       \\
        &          &           & \Machine{L}    \\
        &          &           & \ChoiceC{C}{C} \\[0.8ex]
    \text{terminals}
        & \CTerm{T} & \Coloneqq\ & \CTerm{\Return{V}}     \\
        &           &            & \CTerm{\Lam{x}{M}}
  \end{array}
  \]
\end{minipage}

\bigskip
\raggedright

\fbox{ $\MSteps{L}{C}\phantom{'} $ } \quad \textbf{Basic transitions} \\[1ex]

\emph{Structural transitions}:

\begin{minipage}{0.48\textwidth}
\begin{align*}
  \SLMachine{\CTo{M}{x}{N}}{K}
    &\MStep{} \SLMachine{N}[\AddSusp{S}{x}{M}]{K}
\end{align*}
\end{minipage}
\begin{minipage}{0.48\textwidth}
\begin{align*}
  \SLMachine{\Return{V}}{\StkKont{x}{N}{K}}
    &\MStep{} \SLMachine{\CSubst{N}{V}{x}}{K}
\end{align*}
\end{minipage}

\begin{align*}
  \SLMachine{M}[\DisjSusp{S}{x}{N}]{K}
    &\MStep \SLMachine{N}{\StkKont{x}{M}{K}}
    && \text{(if $x$ is needed in $\CTerm{M}$)} \\[0.5ex]
  \SLMachine{\CApp{M}{V}}{K}
    &\MStep{} \SLMachine{M}{\StkVal{V}{K}} \\
  \SLMachine{\Lam{x}{M}}{\StkVal{V}{K}}
    &\MStep{} \SLMachine{\CSubst{M}{V}{x}}{K} \\
  \SLMachine{\Force*{\Thunk{M}}}{K}
    &\MStep \SLMachine{M}{K} \\
  \SLMachine{\Fix{x}{M}}{K}
    &\MStep \SLMachine{\CSubst{M}{\Thunk*{\Fix{x}{M}}}{x}}{K} \\
  \SLMachine{\CIfz{\Zero}{M}{n}{N}}{K}
    &\MStep \SLMachine{M}{K} \\
  \SLMachine{\CIfz{\Succ{V}}{M}{n}{N}}{K}
    &\MStep \SLMachine{\CSubst{N}{V}{n}}{K} \\
  \SLMachine{\CMatch{\Nil}{M}{x}{xs}{N}}{K}
    &\MStep \SLMachine{M}{K} \\
  \SLMachine{\CMatch{\Cons{V}{W}}{M}{x}{xs}{N}}{K}
    &\MStep \SLMachine{\SubstTwo{N}{V}{W}{x}{xs}}{K}
\end{align*}

\emph{Effect transitions}: \hfill
\begin{minipage}{0.9\textwidth}
\begin{align*}
  \SLMachine{\Fail}{K}
    &\MStep{} \FailC{} \\
  \SLMachine{\Choice{M}{N}}{K}
    &\MStep{} \ChoiceC{\SLMachine{M}{K}}{\SLMachine{N}{K}} \\
  \SLMachine{\Exists{x}[\sigma]{M}}{K}
    &\MStep{} \SLMachine[\Gamma, x : \sigma]{M}{K}
\end{align*}
\end{minipage}

\emph{Narrowing transitions}
  (only if $\SLMachine{M}{K}$ is \emph{blocked} on $x$):
\begin{align*}
  \SLMachine[\Gamma, x : \NatT, \Gamma']{M}{K}
    &\MStep
    \ChoiceC{
      \MSubst{\SLMachine[\Gamma, \Gamma']{M}{K}}{\Zero}{x}
    }{
      \MSubst{\SLMachine[\Gamma, n : \NatT, \Gamma']{M}{K}}{\Succ{n}}{x}
    } \\
  \SLMachine[\Gamma, x : \ListT{\sigma}, \Gamma']{M}{K}
    &\MStep
    \ChoiceC{
      \MSubst{\SLMachine[\Gamma, \Gamma']{M}{K}}{\Nil}{x}
    }{
      \MSubst{\SLMachine[\Gamma, y : \sigma, ys :
      \ListT{\sigma}, \Gamma']{M}{K}}{\Cons{y}{ys}}{x}
    }
\end{align*}

\emph{Unification transitions}:
\begin{align*}
  \SLMachine[\Gamma]{\Equate{V}{W}{M}}{K}
    &\MStep \SimSubM{\prn{\SLMachine[\Gamma]{M}{K}}}{\gamma}
    && (\Mgu{\Gamma,\, \VTerm{V},\, \VTerm{W}} = (\Delta,\, \gamma)) \\
  \SLMachine[\Gamma] {\Equate{V}{W}{M}}{K}
    &\MStep \FailC{}
    && (\Mgu{\Gamma,\, \VTerm{V},\, \VTerm{W}} = \bot)
\end{align*}

\fbox{ $\CSteps{C}{C'}$ } \quad \textbf{Configuration transitions}
\begin{mathpar}
  \inferH{Leaf}{
    \MSteps{L}{C} \\
  }{
    \CSteps{\Ctx{M}[L]}{\Ctx{M}[C]}
  }
  \and
  \inferH{Refl}{
    \strut
  }{
    \CSteps{C}{C}
  }
  \and
  \inferH{Trans}{
    \CSteps{C_1}{C_2} \\
    \CSteps{C_2}{C_3}
  }{
    \CSteps{C_1}{C_3}
  }
\end{mathpar}

  \caption{Abstract Machine. A basic machine $\SLMachine{M}{K}$ carries a
    logic context $\Gamma$, a computation $\CTerm{M}$, a suspension environment
    $\Susp{S}$, and a stack $\Stk{K}$. The bind transition suspends
    $\CTerm{M}$ in $\Susp{S}$ rather than evaluating it immediately; the
    forcing transition reverses this when $x$ is needed. Narrowing transitions
    fire only when a logical variable blocks the current computation.}
  \label{figure:machine}
\end{figure*}

The CBPV term produced by the translation of \cref{section:ir} is not evaluated
directly. Instead, it gives rise to a \emph{nondeterministic search tree}: a
(possibly infinite) tree whose internal nodes are choice points and whose leaves
are either results or failures. The task of execution is to explore this tree
and collect the results. We decompose this task into two parts: a
\emph{basic machine} that steps a single CBPV computation until it reaches a
branch point, and a \emph{search strategy} that manages the collection of
active machines.

This decomposition is a direct realisation of the abstract machine of
\parencite{jones_2026}. That paper organises machine states into
\emph{configurations}: binary trees whose leaves are basic machines or
failures. A configuration transition selects a non-terminal leaf and steps it,
possibly splitting it into further branches. The theory is silent on
\emph{which} leaf to step:
%
\begin{quote}
  \emph{``Observe that this is highly nondeterministic: it may be that one or
  more leaves of $\Conf{C}$ may be able to make a basic transition! We leave
  the choice to the implementer, who may use any strategy to explore the tree
  (e.g.~DFS, BFS, parallel).''} \hfill---\parencite{jones_2026}
\end{quote}
%
\noindent This freedom is justified by the denotational semantics: because the
semantics is based on the lower powerdomain, which is insensitive to the order
of exploration, all strategies produce the same set of results for terminating
programs.

\subsection{The Basic Machine}

The basic machine (\cref{figure:machine}) is a state machine whose transitions
correspond to the evaluation of a single CBPV computation. Its state consists
of:
%
\begin{itemize}
  \item A \emph{computation closure} $(\CTerm{M}, \rho)$: the current
    computation paired with an environment $\rho$ mapping de Bruijn indices to
    value closures.
  \item A \emph{stack} $\Stk{K}$: a list of frames representing pending
    continuations ($\textsf{to}\ x.\, \CTerm{N}$), argument values
    ($\VTerm{V}$), and suspension-set frames that record which suspension to
    update when a forced computation completes.
  \item A \emph{logic context} $\Gamma$: a record of logical variables, their
    first-order types, and their current bindings.
  \item A \emph{suspension environment} $\Susp{S}$: a map from suspension
    identifiers to deferred computation closures (discussed in
    \cref{section:suspensions}).
\end{itemize}

\noindent The machine steps by pattern-matching on the current computation. Most
transitions are standard CBPV reductions: a $\CTerm{\Return{V}}$ pops a
continuation from the stack and substitutes $\VTerm{V}$ into the continuation
body; an application $\CTerm{\CApp{M}{V}}$ pushes $\VTerm{V}$ onto the stack;
forcing a thunk $\CTerm{\Force*{\Thunk{M}}}$ unwraps the suspended computation;
a case expression selects the appropriate branch based on the head constructor
of the scrutinee. Each of these transitions corresponds to a $\beta$-law of the
equational theory (\cref{figure:equations}).

The machine is structured around a tight inner loop that executes deterministic
steps until a branch point (choice or narrowing) or a terminal state (return
with empty stack, or failure) is reached. Control is returned to the search
strategy only at points where a genuine scheduling decision must be made; all
deterministic transitions are executed in place.

\subsection{Effect Transitions}

The FLP effects produce the branching structure of the search tree.

\paragraph{Choice}
When the machine encounters $\CTerm{\Choice{M}{N}}$, it produces a
\emph{branch point}: two copies of the current state, one continuing with
$\CTerm{M}$ and the other with $\CTerm{N}$. This is the fundamental source of
nondeterminism. The equational theory guarantees that choice is associative,
commutative, and idempotent, so the tree's branching structure does not affect
the set of results.

\paragraph{Failure}
$\CTerm{\Fail}$ discards the current branch: it is a leaf of the search tree
with no result.

\paragraph{Logical variables}
$\CTerm{\Exists{x}[\sigma]{M}}$ allocates a fresh entry in the logic context
with type $\sigma$ and no binding, extends the environment with a reference to
this entry, and continues with $\CTerm{M}$. The variable $x$ is initially
unconstrained; its value is determined later by narrowing or unification.

\paragraph{Narrowing}
The critical interaction between logical variables and case analysis is handled
by narrowing. When the machine encounters a case expression whose scrutinee
resolves---through the environment and logic context---to an \emph{unbound}
logical variable $x : \sigma$, it cannot proceed: the variable's head
constructor is unknown. The machine responds by \emph{splitting} into branches,
one for each constructor of $\sigma$. In each branch, $x$ is bound to the
appropriate constructor applied to fresh logical variables.

For example, if $x : \NatT$ is unbound and the machine reaches
$\textsf{ifz}(x; \ldots)$, it produces two branches: one binding $x$ to
$\VTerm{\Zero}$, and one binding $x$ to $\VTerm{\Succ{y}}$ for a fresh
$y : \NatT$. This corresponds to the \emph{parameter equation}
$\CTerm{\Exists{n}[\NatT]{M}} =
\CTerm{\Choice{\CSubst{M}{\Zero}{n}}{\Exists{m}[\NatT]{\CSubst{M}{\Succ{m}}{n}}}}$
of \cref{figure:equations}: the equation says that introducing a logical
variable of type $\NatT$ is the same as choosing between zero and the successor
of a fresh variable.

The machine narrows only when a variable \emph{blocks} a computation, in the
sense of \parencite{jones_2026}: the next step directly depends on the
variable's value. This strategy, called \emph{weak head narrowing}, ensures that
narrowing happens only when necessary. Its completeness---every closed value of
type $\sigma$ is eventually considered as a possible instantiation---is proved
in \parencite{jones_2026}.

\paragraph{Unification}
$\CTerm{\Equate{V}{W}{M}}$ attempts to unify $\VTerm{V}$ and $\VTerm{W}$ using
the standard first-order unification algorithm. The algorithm processes a queue
of pairs of value closures, decomposing constructors and binding variables via
the logic context. Three outcomes are possible: \emph{success}, in which the
logic context is updated with the most general unifier and execution continues
with $\CTerm{M}$; \emph{failure}, in which the values have incompatible head
constructors and the branch is discarded; or \emph{suspension}, in which one of
the values refers to a deferred computation that must be forced before
unification can proceed.

An optional occurs check prevents the construction of infinite terms (e.g.\
$x = \textsf{S}(x)$). It can be disabled for performance at the cost of
soundness.

\subsection{Search Strategies}
\label{subsection:search-strategies}

The search tree generated by the machine must be explored to collect results.
\Gweek{} implements four strategies, each corresponding to a different
traversal of the tree.

\paragraph{Breadth-first search (\BFS{})}
The default strategy maintains two vectors: the current generation and the next
generation. Each machine in the current generation is stepped to its next branch
point; the resulting branches are placed in the next generation. When the
current generation is exhausted, the vectors are swapped. \BFS{} is
\emph{complete}---if a solution exists at any finite depth, it will be
found---and \emph{fair}: no branch is starved. Its limitation is memory: the
number of active machines can grow exponentially in the depth of the tree.

\paragraph{Depth-first search (\DFS{})}
\DFS{} maintains a stack and always steps the most recently created branch. It
uses memory proportional to the depth of the tree, but it is
\emph{incomplete}: if the tree has an infinite branch, \DFS{} may follow it
forever. Infinite branches arise naturally when logical variables of type
$\NatT$ or $\ListT{A}$ are narrowed, so \DFS{} is appropriate only for finite
search spaces or when a single solution is sufficient.

\paragraph{Iterative deepening (\IDDFS{})}
Iterative deepening combines the completeness of \BFS{} with the memory
efficiency of \DFS{}. The strategy runs \DFS{} with a depth limit that starts
at~1 and doubles after each complete pass. Solutions found in previous passes
are deduplicated by a hash set of previously seen outputs. The overhead of
re-exploring shallow nodes is bounded by a constant factor.

\paragraph{Fair round-robin (\Fair{})}
The fair strategy maintains a queue of \emph{threads}, where each thread is a
local \DFS{} work stack. A single thread is created initially; new threads are
created at branch points. When the machine encounters a branch while executing a
thread, the first alternative continues in the current thread and each remaining
alternative is placed in the queue as a new single-machine thread. Each thread
is allowed a fixed quota of steps (currently $10{,}000$). If a thread exhausts
its quota without completing, its remaining work stack is re-enqueued as a
single entry, and the scheduler moves to the next thread.

The key invariant is that \emph{branch alternatives are distributed across
threads}, so every branch receives its own share of the round-robin schedule.
Within a thread, execution is depth-first, so the local work stack remains
small (proportional to the depth of the sub-tree being explored). To bound
memory, the total number of threads is capped at a configurable maximum
(currently $10{,}000$); when the cap is reached, new branch alternatives remain
in their parent thread and are explored depth-first. This degrades gracefully:
the branches that have already been distributed across threads continue to
receive fair scheduling, while only new branches within a full thread lose
individual time-slicing.

This strategy is both complete and fair: every branch is eventually explored,
and no branch is starved indefinitely. It avoids the memory blow-up of \BFS{}
while retaining completeness, and is particularly effective for programs that
mix finite branching (explicit choice) with infinite branching (logical variable
narrowing).

\paragraph{Discussion}
The theoretical licence to choose any exploration order is powerful, but it
places a practical burden on the programmer. In general, \BFS{} is the safest
default for programs involving logical variables, as it does not require the
programmer to reason about the finiteness of branches. \DFS{} is fastest for
finite search spaces. The fair strategy offers a practical compromise when the
structure of the search space is unknown.


\section{The Rust Implementation}
  \label{section:implementation}
  The implementation of \Gweek{} is written in Rust and consists of approximately
3{,}000 lines of code. In this section we describe the compilation pipeline and
the key engineering decisions that make the implementation practical.

\subsection{Pipeline}

A \Gweek{} program passes through four stages before execution. The source text
is first \emph{parsed} into an AST of declarations and expressions, using
parser combinators (the \texttt{chumsky} crate). Each top-level declaration
consists of a type signature and a definition; the body is a single expression
in the syntax of \cref{section:source}. The AST is then \emph{type-checked}
against the type system of \cref{figure:cbv}: the checker infers types for
local bindings and verifies that existential quantifiers range only over
first-order types.

The type-checked AST is then \emph{translated} into CBPV machine terms,
following the translation of \cref{subsection:translation}. Values become
\texttt{MValue} terms and computations become \texttt{MComputation} terms.
Named variables are replaced by de Bruijn indices: each variable occurrence
becomes an index into an environment of value closures, eliminating the need for
$\alpha$-renaming during substitution. Function definitions are topologically
sorted and compiled into an environment of closures; each definition becomes a
fixed-point $\CTerm{\Fix{f}{\Lam{x}{M}}}$ over a thunked computation. An
optional peephole optimiser (\cref{section:optimisation}) is then applied to the
resulting terms.

Finally, the translated term is handed to the abstract machine of
\cref{section:machine}, which generates a nondeterministic search tree. A search
strategy (\cref{subsection:search-strategies}) explores the tree and collects
results, printing them as they are discovered.

\subsection{Arena Allocation}

All values, computations, environments, and stack frames are allocated in a
single bump arena (using the \texttt{bumpalo} crate). A bump arena allocates by
advancing a pointer, making allocation essentially free; deallocation is
all-or-nothing, with the entire arena freed when the program terminates.

This design is well-suited to the execution model. The machine continuously
creates new terms, environments, and stack frames as it steps, but it never
deallocates individual objects during execution. The arena also enables a
simple representation of sharing: closures and environments are arena-allocated
references that can be freely duplicated without the overhead of reference
counting or garbage collection.

\subsection{Persistent Data Structures and Cheap Branching}

The central performance requirement of an FLP implementation is that
\emph{branching must be cheap}. At every choice point or narrowing split, the
machine duplicates its state so that each branch can proceed independently. If
this duplication were $O(n)$ in the size of the environment or stack, the cost
of nondeterminism would quickly become prohibitive.

We address this by representing the variable environment and the stack as
\emph{persistent cons-lists} allocated in the arena. Each node is immutable;
extending an environment or pushing a stack frame allocates a single new cons
cell that points to the existing list. Duplicating a machine state therefore
copies only the root pointers---an $O(1)$ operation---and the two branches
share the entire prefix of their environments and stacks.

The logic context, which tracks logical variables and their bindings, uses a
different strategy. It is represented as a reference-counted vector of entries
(each storing a type and an optional value binding) paired with a union-find
structure for variable aliasing. When the machine branches, both children
initially share the same backing storage. If either branch subsequently modifies
the context---by binding a variable during narrowing or unification---the
storage is copied on write via \texttt{Rc::make\_mut}. This \emph{copy-on-write}
discipline ensures that branching is cheap in the common case (no copy is made
until a modification occurs) while maintaining the invariant that each branch
has an independent view of the logic context. The suspension environment uses
the same copy-on-write strategy: a reference-counted vector of entries, each
either a pending computation closure or a resolved value closure.

\subsection{Value Closures}

Because the machine uses de Bruijn indices, a value in isolation is not
meaningful: it must be interpreted relative to an environment $\rho$ that maps
its free indices to actual values. We call the pair of a value and its
environment a \emph{value closure}. Value closures come in three variants:
\texttt{Clos}$(\VTerm{V}, \rho)$ is an ordinary value paired with its
environment; \texttt{LogicVar}$(x)$ is a reference to a logical variable in the
logic context; and \texttt{Susp}$(s)$ is a reference to a suspension in the
suspension environment.

The operation \texttt{close\_head} resolves a value closure to its \emph{head
form}: the outermost constructor or an unbound logical variable. It works by
following a chain of indirections: if the closure is a \texttt{Clos} wrapping a
de Bruijn variable, it looks up the index in the environment; if it is a
\texttt{LogicVar} with a binding in the logic context, it follows the binding;
if it is a \texttt{Susp}, it checks whether the suspension has been resolved.
This loop terminates when it reaches a constructor, an unbound logical variable,
or an unresolved suspension. In the last case, \texttt{close\_head} does not
transparently force the suspension; instead, it signals to the caller that
forcing is required, and the machine initiates the forcing protocol described
in \cref{section:suspensions}.

This operation is called at every case expression and every unification step,
and must be efficient. It terminates because the chain of indirections is always
finite---an invariant maintained by the occurs check in unification.


\section{Optimisation}
  \label{section:optimisation}
  A distinctive feature of \Gweek{} is that its optimiser is derived directly from
the equational theory. The equations of \cref{figure:equations}, validated by
the denotational semantics of \parencite{jones_2026}, serve as rewrite rules:
each rule replaces a term with an equal one, and the soundness of the
transformation follows from the full abstraction theorem of the semantics rather
than from ad hoc argument.

The optimiser operates as a \emph{peephole pass} over the CBPV term produced by
the translation of \cref{subsection:translation}. It traverses the term
bottom-up, applying rewrites at each node until a fixed point is reached.
Because the machine uses de Bruijn indices, the rewrites require shifting and
substitution on indices, which we implement as straightforward recursive
traversals. \Cref{figure:rewrites} lists all rewrite rules alongside the
equations that justify them. We now describe them in detail.

\begin{figure}
\begin{small}
\begin{tabular}{@{}lll@{}}
  \toprule
  \textbf{Rewrite} & \textbf{Equation} & \textbf{Law} \\
  \midrule
  \textsc{bind-eta}
    & $\CTerm{\CTo{\Return{V}}{x}{M}} \leadsto \CTerm{\CSubst{M}{V}{x}}$
    & $\beta$-bind \\
  \textsc{dead-bind}
    & $\CTerm{\CTo{\Return{V}}{x}{N}} \leadsto \CTerm{N}$ \quad ($x \notin \textsf{fv}(N)$)
    & $\beta$-bind \\
  \textsc{bind-assoc}
    & $\CTerm{\CTo*{\CTo{M}{x}{N}}{y}{P}} \leadsto \CTerm{\CTo{M}{x}{(\CTo{N}{y}{P})}}$
    & seq.\ assoc. \\
  \textsc{app-bind}
    & $\CTerm{\CApp*{\CTo{M}{x}{N}}{V}} \leadsto \CTerm{\CTo{M}{x}{\CApp{N}{V}}}$
    & seq.\ app. \\
  \textsc{lam-beta}
    & $\CTerm{\CApp*{\Lam{x}{M}}{V}} \leadsto \CTerm{\CSubst{M}{V}{x}}$
    & $\beta$-app \\
  \textsc{force-beta}
    & $\CTerm{\Force*{\Thunk{M}}} \leadsto \CTerm{M}$
    & $\beta$-thunk \\
  \textsc{ifz-beta}
    & $\CTerm{\CIfz{\Zero}{M}{n}{N}} \leadsto \CTerm{M}$ (etc.)
    & $\beta$-ifz \\
  \textsc{match-beta}
    & $\CTerm{\CMatch{\Cons{V}{W}}{M}{x}{xs}{N}} \leadsto \CTerm{\SubstTwo{N}{V}{W}{x}{xs}}$ (etc.)
    & $\beta$-match \\
  \textsc{case-beta}
    & $\CTerm{\CCase{\In[1]{V}}{x}{M}{y}{N}} \leadsto \CTerm{\CSubst{M}{V}{x}}$ (etc.)
    & $\beta$-case \\
  \textsc{dead-end}
    & $\CTerm{\CTo{M}{x}{\Fail}} \leadsto \CTerm{\Fail}$ \quad ($M$ total)
    & dead-end \\
  \textsc{lam-choice}
    & $\CTerm{\Lam*{x}{\Choice{M}{N}}} \leadsto \CTerm{\Choice{\Lam{x}{M}}{\Lam{x}{N}}}$
    & $\lambda$-$\talloblong$ \\
  \textsc{lam-exists}
    & $\CTerm{\Lam*{x}{\Exists{z}[\sigma]{M}}} \leadsto \CTerm{\Exists*{z}[\sigma]{\Lam{x}{M}}}$
    & $\lambda$-$\exists$ \\
  \textsc{lam-equate}
    & $\CTerm{\Lam*{x}{\Equate{V}{W}{M}}} \leadsto \CTerm{\Equate{V}{W}{\Lam{x}{M}}}$
    & $\lambda$-eq \\
  \textsc{eq-choice}
    & $\CTerm{\Equate{V}{W}{\Choice{M}{N}}} \leadsto \CTerm{\Choice{(\Equate{V}{W}{M})}{(\Equate{V}{W}{N})}}$
    & eq-$\talloblong$ \\
  \textsc{eq-exists}
    & $\CTerm{\Equate{V}{W}{\Exists{z}[\sigma]{M}}} \leadsto \CTerm{\Exists*{z}[\sigma]{\Equate{V}{W}{M}}}$
    & eq-$\exists$ \\
  \textsc{cycle}
    & $\CTerm{\Equate{V}{\Succ{V}}{M}} \leadsto \CTerm{\Fail}$ (etc.)
    & occurs \\
  \bottomrule
\end{tabular}
\end{small}
\caption{Peephole rewrite rules and their justifying equations. All equations
  appear in \cref{figure:equations} or are validated by the denotational
  semantics of \parencite{jones_2026}. In addition, the choice identity laws
  ($\CTerm{\Choice{\Fail}{M}} = \CTerm{M}$), the idempotence law
  ($\CTerm{\Choice{M}{M}} = \CTerm{M}$), the reflexivity of unification
  ($\CTerm{\Equate{V}{V}{M}} = \CTerm{M}$), and the narrowing
  decomposition rules are applied.}
\label{figure:rewrites}
\end{figure}

\paragraph{$\beta$-reductions}
The $\beta$-laws of \cref{figure:equations} are the workhorses of the
optimiser. The bind-return law
$\CTerm{\CTo{\Return{V}}{x}{M}} = \CTerm{\CSubst{M}{V}{x}}$
(\textsc{bind-eta}) is the most frequently applied: it eliminates the overhead
of a continuation push and pop by substituting the value directly. A special
case arises when the returned value is a variable
($\CTerm{\CTo{\Return{x}}{y}{M}}$), which is reduced by aliasing $y$ to $x$.
When $x$ does not appear free in the continuation, the bind is dead and the
continuation is returned unchanged (\textsc{dead-bind}).

The remaining $\beta$-laws reduce eliminators applied to known constructors:
\textsc{ifz-beta} reduces $\CTerm{\CIfz{\Zero}{M}{n}{N}}$ to $\CTerm{M}$ and
$\CTerm{\CIfz{\Succ{V}}{M}{n}{N}}$ to $\CTerm{\CSubst{N}{V}{n}}$;
\textsc{match-beta} reduces
$\CTerm{\CMatch{\Cons{V}{W}}{M}{x}{xs}{N}}$ to
$\CTerm{\SubstTwo{N}{V}{W}{x}{xs}}$;
\textsc{case-beta} reduces
$\CTerm{\CCase{\In[1]{V}}{x}{M}{y}{N}}$ to $\CTerm{\CSubst{M}{V}{x}}$.
The \textsc{force-beta} rule reduces $\CTerm{\Force*{\Thunk{M}}}$ to
$\CTerm{M}$, and \textsc{lam-beta} reduces
$\CTerm{\CApp*{\Lam{x}{M}}{V}}$ to $\CTerm{\CSubst{M}{V}{x}}$.

\paragraph{Sequencing laws}
The associativity law (\textsc{bind-assoc}) reassociates
$\CTerm{\CTo*{\CTo{M}{x}{N}}{y}{P}}$ to
$\CTerm{\CTo{M}{x}{(\CTo{N}{y}{P})}}$, exposing opportunities for
\textsc{bind-eta}. The application-bind law (\textsc{app-bind}) rewrites
$\CTerm{\CApp*{\CTo{M}{x}{N}}{V}}$ to $\CTerm{\CTo{M}{x}{\CApp{N}{V}}}$.

\paragraph{$\lambda$-laws}
The optimiser distributes $\lambda$-abstractions over effects using the three
$\lambda$-laws of \cref{figure:equations}: \textsc{lam-choice} rewrites
$\CTerm{\Lam*{x}{\Choice{M}{N}}}$ to
$\CTerm{\Choice{\Lam{x}{M}}{\Lam{x}{N}}}$; \textsc{lam-exists} rewrites
$\CTerm{\Lam*{x}{\Exists{z}[\sigma]{M}}}$ to
$\CTerm{\Exists*{z}[\sigma]{\Lam{x}{M}}}$; and \textsc{lam-equate} rewrites
$\CTerm{\Lam*{x}{\Equate{V}{W}{M}}}$ to
$\CTerm{\Equate{V}{W}{\Lam{x}{M}}}$ (when $\VTerm{V}$ and $\VTerm{W}$ do not
reference $x$). These normalise terms so that effects appear at the outermost
level.

\paragraph{Effect equations}
The identity law $\CTerm{\Choice{\Fail}{M}} = \CTerm{M}$ removes dead branches,
and the idempotence law $\CTerm{\Choice{M}{M}} = \CTerm{M}$ collapses
duplicates. Nested choices are flattened using associativity. The dead-end law
(\textsc{dead-end}) rewrites $\CTerm{\CTo{M}{x}{\Fail}}$ to $\CTerm{\Fail}$
when $\CTerm{M}$ is total---a condition we check by a conservative syntactic
analysis. Reflexive unification $\CTerm{\Equate{V}{V}{M}} = \CTerm{M}$
eliminates trivially satisfied constraints.

\paragraph{Equational constraint rules}
The \textsc{eq-choice} rule distributes an equational constraint over a choice:
$\CTerm{\Equate{V}{W}{\Choice{M}{N}}} =
\CTerm{\Choice{(\Equate{V}{W}{M})}{(\Equate{V}{W}{N})}}$.
The \textsc{eq-exists} rule commutes an equational constraint past an
existential: $\CTerm{\Equate{V}{W}{\Exists{z}[\sigma]{M}}} =
\CTerm{\Exists*{z}[\sigma]{\Equate{V}{W}{M}}}$.
The narrowing equations decompose unification on known constructors: for
example, $\CTerm{\Equate{\Succ{V}}{\Succ{W}}{M}}$ reduces to
$\CTerm{\Equate{V}{W}{M}}$, and
$\CTerm{\Equate{\Cons{V_1}{W_1}}{\Cons{V_2}{W_2}}{M}}$ reduces to
$\CTerm{\Equate{V_1}{V_2}{\Equate{W_1}{W_2}{M}}}$.
Mismatched constructors ($\CTerm{\Equate{\Succ{V}}{\Zero}{M}}$,
$\CTerm{\Equate{\Cons{V}{W}}{\Nil}{M}}$, etc.)\ reduce to $\CTerm{\Fail}$.
The \textsc{cycle} rule detects cyclic constraints
($\CTerm{\Equate{V}{\Succ{V}}{M}}$, $\CTerm{\Equate{W}{\Cons{V}{W}}{M}}$,
etc.)\ and reduces them to $\CTerm{\Fail}$.

\paragraph{Discussion}
The key point is not the magnitude of the improvement but its
\emph{trustworthiness}: because every rewrite is a validated equation, the
optimiser cannot change the observable behaviour of the program. This is a
stronger guarantee than is typical for compiler optimisations, and it is a
direct consequence of basing the implementation on the algebraic theory.


\section{Suspensions}
  \label{section:suspensions}
  A distinctive consequence of basing \Gweek{} on CBPV is its treatment of
\texttt{let} bindings. In the core calculus, a \texttt{let} binding corresponds
to the sequencing construct $\CTerm{\CTo{M}{x}{N}}$---the monadic bind. In a
strict evaluator, $\CTerm{M}$ would be fully evaluated before control passes to
$\CTerm{N}$. But in a language with nondeterminism, this is not the only
reasonable choice, and the difference turns out to be consequential.

\subsection{The Mechanism}

The key transition is the bind rule in \cref{figure:machine}. In a standard CK
machine, $\CTerm{\CTo{M}{x}{N}}$ pushes a continuation onto the stack and
evaluates $\CTerm{M}$ immediately. \Gweek{} instead \emph{suspends}
$\CTerm{M}$: it records the computation in the suspension environment
$\Susp{S}$ and continues with $\CTerm{N}$:
%
\[
  \SLMachine{\CTo{M}{x}{N}}{K}
  \quad\MStep\quad
  \SLMachine{N}[\AddSusp{S}{x}{M}]{K}
\]
%
The computation $\CTerm{M}$ is not evaluated. The suspension environment is
extended with the binding $x \mapsto \CTerm{M}$, and the variable $x$ in
$\CTerm{N}$ refers to this suspended computation.

The suspension is \emph{forced}---the deferred computation is actually
evaluated---only when $x$ is encountered in a position that \emph{demands} its
value: a case expression whose scrutinee is $x$, or a unification whose operand
is $x$. When this happens, the machine reverses the suspension: it extracts the
suspended computation from $\Susp{S}$, pushes a special \emph{set} frame onto
the stack that records which suspension to update, and begins evaluating the
deferred computation:
%
\[
  \SLMachine{M}[\DisjSusp{S}{x}{N}]{K}
  \quad\MStep\quad
  \SLMachine{N}{\StkKont{x}{M}{K}}
  \qquad\text{(when $x$ is needed in $\CTerm{M}$)}
\]
%
When the deferred computation completes with a value, the set frame updates the
suspension environment so that subsequent references to $x$ resolve immediately.
The effect is exactly what the strict machine would have done---but deferred
until the point of use.

If $x$ is never needed during normal execution---if it is used only in
positions that do not inspect its value, such as being passed as an argument or
stored in a data structure---the suspension remains unforced until the machine
reaches a terminal state (a return with an empty stack). At that point the
machine enters a \emph{drain phase}: it iterates over the remaining pending
suspensions, forcing each in turn by pushing a set frame and evaluating the
deferred computation. This ensures that any effects the suspensions contain (in
particular, failures that should prune the current branch) are properly
accounted for before a result is reported.

\subsection{Why Suspensions Matter}

The difference between the strict and lazy transitions may seem minor, but it
has a profound effect on the interaction between sequencing and nondeterminism.

Consider a program that searches for a sorted permutation of a list:
%
\begin{haskellcode}
let ys = perm [3, 1, 2] in
let c = guard_sorted ys in
ys.
\end{haskellcode}
%
\noindent Under the strict semantics, the machine evaluates
\mintinline{haskell}|perm [3, 1, 2]| first. This is a nondeterministic
computation that branches into all $3! = 6$ permutations. Each branch is a
separate machine state. Only then does the machine evaluate
\mintinline{haskell}|guard_sorted ys| in each branch independently. The
constraint check cannot prune branches that were created \emph{before} it ran.

Under the suspension semantics, the evaluation proceeds differently. First,
\mintinline{haskell}|perm [3, 1, 2]| is suspended---no branches are created
yet. Then \mintinline{haskell}|guard_sorted ys| is also suspended. The machine
reaches the result expression \mintinline{haskell}|ys|, which is a suspension
reference. If the context demands the value of \mintinline{haskell}|ys|---for
example, to print the result---then \mintinline{haskell}|ys| is forced. The
permutation is generated incrementally: each element is produced and immediately
available to the sortedness check. When \mintinline{haskell}|guard_sorted|
examines the first two elements and finds them unsorted, it fails, pruning the
branch \emph{before} the remaining elements are ever chosen.

The equational theory illuminates this. The
\emph{bind-return law} states that $\CTerm{\CTo{\Return{V}}{x}{M}} =
\CTerm{\CSubst{M}{V}{x}}$: when the bound computation simply returns a value,
sequencing is equivalent to substitution. But when $\CTerm{M}$ is
nondeterministic---when it branches into multiple alternatives---the bind
genuinely sequences, and the continuation $\CTerm{N}$ runs in each branch. The
suspension mechanism defers this branching, allowing the continuation to proceed
and potentially constrain the search before the full tree is materialised.

\subsection{Strict Mode}

\Gweek{} supports a strict evaluation mode (the \texttt{--strict} flag) in
which every \texttt{let} binding uses the standard CK transition instead of the
suspension transition. The machine pushes a continuation frame and evaluates the
bound computation immediately.

Strict mode is appropriate for programs whose search space is inherently
finite, or where all constraints are forced immediately by case expressions. It
eliminates the overhead of the suspension environment and avoids the subtleties
of deferred evaluation.

However, strict mode changes the termination behaviour of the system. Programs
that rely on lazy interleaving of generation and testing may fail to terminate
in strict mode, because the generation fully unfolds before any constraint can
prune it. The spectrum between lazy and strict FLP semantics---and the question
of which equations hold in each---is an interesting direction for future work.

\subsection{Programming Discipline}

The suspension mechanism requires the programmer to attend to two patterns.

\paragraph{Forcing constraints}
If a constraint is bound to a variable that is never scrutinised, the
constraint is suspended and never forced during execution. For example, in
\mintinline{haskell}|let c = neq a b in body|, if \mintinline{haskell}|c| is
never case-analysed in \mintinline{haskell}|body|, the call to
\mintinline{haskell}|neq| has no effect. The fix is to force the result:
%
\begin{haskellcode}
case neq a b of
    Z -> body
  | S n -> fail.
\end{haskellcode}
%
\noindent The case scrutinee demands the head constructor, which forces the
suspension and triggers any embedded failure.

\paragraph{Interleaving generation and testing}
Even when constraints are properly forced, the order in which nondeterministic
choices and constraints are combined matters for efficiency. Generating all
candidates monolithically and then testing each means the full nondeterministic
tree is explored before any constraint can prune it. The efficient alternative
is to \emph{interleave} generation and testing: produce one candidate element,
check applicable constraints, and only then produce the next.


\section{Benchmarking}
  \label{section:benchmarking}
  \input{benchmarking}

\section{Related Work}
  \label{section:related-work}
  \paragraph{Functional Logic Programming Languages}

The combination of functional and logic programming has a long history, surveyed
by
\textcite{bellia_1986,darlington_1986,hanus_1994,antoy_2010}. The most
prominent FLP language is Curry \parencite{hanus_2013,hanus_2016}, which extends
Haskell with logical variables, equational constraints, and nondeterministic
functions. Curry has multiple implementations, including KiCS2
\parencite{brassel_2011}, which compiles Curry to Haskell, and PAKCS
\parencite{hanus_2022}, which targets Prolog. \Gweek{} is intentionally much
simpler than Curry---it lacks type classes, modules, and I/O---but it makes the
underlying algebraic theory explicit in its design.

The Verse language \parencite{EpicGames_2023,augustsson_2023} is a recent FLP
language whose core calculus is an untyped $\lambda$-calculus with choice and
logical variables. Verse uses a stream semantics and supports encapsulated
search. \Gweek{} differs from Verse in being typed, in using CBPV as its
foundation, and in not supporting encapsulated search. As discussed in
\parencite{jones_2026}, encapsulated search requires moving beyond algebraic
effects to \emph{effect handlers}, and raises questions about the interaction of
encapsulation with the commutativity and idempotence of choice.

The miniKanren family of relational programming languages
\parencite{friedman_2005,friedman_2018,byrd_2009} embeds logic programming in a
host language (typically Scheme). miniKanren is conceptually close to \Gweek{} in
its use of logical variables and unification, but it is purely relational: it
lacks the functional programming features (higher-order functions, algebraic
data types, pattern matching) that \Gweek{} inherits from CBPV.

\paragraph{Semantics of FLP}

The semantic foundations of \Gweek{} are developed in our companion paper
\parencite{jones_2026}. That paper presents a domain-theoretic semantics for a
CBPV-based FLP calculus, based on the lower powerdomain, and proves it sound,
adequate, and fully abstract for a termination-insensitive notion of
observational equivalence. The equational theory derived from this semantics
validates the rewrite rules used in \Gweek{}'s peephole optimiser.

Previous semantic work on FLP includes the CRWL rewriting logic framework of
\textcite{gonzalez-moreno_1999}, the operational semantics of
\textcite{albert_2005}, and the denotational approaches of
\textcite{christiansen_2011} and \textcite{mehner_2014}. The monad-based
approach to FLP effects, which \Gweek{} follows, originates with
\textcite{staton_2013}. The observation that logical variables are an effect that
can be added to a language in isolation goes back to \textcite{lindstrom_1985}.

\paragraph{Narrowing}

Narrowing \parencite{reddy_1985} is the standard execution strategy for FLP.
\Gweek{} implements \emph{weak head narrowing}: a logical variable is refined
only when it is scrutinised by a case expression or required by unification.
This is related to the \emph{needed narrowing} strategy of
\textcite{antoy_1997}, which refines variables only when they are `needed' by
the outermost function. In the CBPV framework, the distinction between values
and computations provides a clean account of when narrowing must occur: it
happens precisely when a computation attempts to eliminate a value that turns out
to be a logic variable. The completeness of this strategy is proved in
\parencite{jones_2026}.

\paragraph{Implementation Techniques}

\Gweek{}'s use of arena allocation for term representation is standard in
high-performance language implementations; see e.g.\ the Zig compiler and the
Rust compiler's use of arenas. The persistent cons-list representation of
environments is a classical technique in functional language implementations
\parencite{landin_1964,henderson_1980}. The union-find structure for logic
variables is the standard choice in logic programming implementations
\parencite{warren_1983}.

Embeddings of logic programming in functional languages
\parencite{claessen_2001,seres_1999,hinze_1998} provide similar functionality
to \Gweek{} but at the cost of embedding overhead and the loss of the host
language's type system for the logical fragment. As a standalone language,
\Gweek{} retains control over term representation and search implementation.


\section{Conclusions}
  \label{section:conclusions}
  We have presented \Gweek{}, an implementation of a functional logic programming
language whose design is derived from the algebraic theory of FLP effects. The
key insight driving the design is that the characteristic constructs of
FLP---nondeterministic choice, logical variables, and equational
constraints---are algebraic effects in the sense of \textcite{plotkin_2003}, and
can be given a clean account in the CBPV framework of \textcite{levy_2003}.

The implementation demonstrates that the algebraic theory is not merely a
theoretical curiosity: it provides direct guidance for the design of the
abstract machine, the equational laws used by the optimiser, and the
understanding of when and how effects interact. The suspension mechanism, in
particular, arises naturally from the CBPV treatment of sequencing, and its
consequences for programming---the need to force constraints, the interaction
between laziness and nondeterminism, the dramatic effect of interleaving
generation and testing---are illuminated by the equational theory.

Several directions for future work present themselves. First, the language could
be extended with \emph{encapsulated search}: a primitive that collects all
results of a nondeterministic computation into a value. As discussed in
\parencite{jones_2026}, this requires moving beyond algebraic effects to
\emph{effect handlers}, and raises interesting questions about the interaction of
encapsulation with the commutativity and idempotence of choice. Second, the type
system could be extended with polymorphism and type classes, enabling more
expressive programs. Third, the implementation could benefit from more
sophisticated search strategies, such as constraint propagation or learned
heuristics, which would bring \Gweek{} closer to the state of the art in
constraint programming. Finally, the strict evaluation mode suggests an
interesting spectrum between lazy and strict FLP semantics, whose theoretical
and practical implications deserve further study.


\printbibliography

\end{document}
